\section*{Chapitre \ref{ch:g9:fractions} --- Fractions et puissances}

\begin{solution}{exo:g9:fractions:1}
$\dfrac13 + \dfrac15 = \dfrac{5}{15} + \dfrac{3}{15} = \dfrac{8}{15}$.

$\dfrac72 - \dfrac53 = \dfrac{21}{6} - \dfrac{10}{6} = \dfrac{11}{6}$.

$\dfrac{5}{12} + \dfrac38 = \dfrac{10}{24} + \dfrac{9}{24} =
\dfrac{19}{24}$.

$3 - \dfrac45 = \dfrac{15}{5} - \dfrac45 = \dfrac{11}{5}$.
\end{solution}

\begin{solution}{exo:g9:fractions:2}
$\dfrac59 \times \dfrac{27}{35} = \dfrac{5 \times 27}{9 \times 35}
= \dfrac{27}{9} \times \dfrac{5}{35} = 3 \times \dfrac17 = \dfrac37$
(simplifier par $9$ et par $5$).

$\dfrac{8}{15} \div \dfrac45 = \dfrac{8}{15} \times \dfrac54
= \dfrac{8 \times 5}{15 \times 4} = \dfrac{40}{60} = \dfrac23$.

$\dfrac23 \times \dfrac34 \times \dfrac45 = \dfrac{2 \times 3 \times 4}
{3 \times 4 \times 5} = \dfrac25$ (les $3$ et les $4$ s'annulent).
\end{solution}

\begin{solution}{exo:g9:fractions:3}
$2^5 = 32$~; \quad $(-2)^4 = 16$~; \quad $-2^4 = -16$~; \quad
$4^{-2} = \frac{1}{16}$~; \quad
$\left(\frac23\right)^3 = \frac{8}{27}$~; \quad $7^0 = 1$.
\end{solution}

\begin{solution}{exo:g9:fractions:4}
$5^3 \times 5^6 = 5^9$~; \quad $\dfrac{7^9}{7^5} = 7^4$~; \quad
$\left(2^3\right)^4 = 2^{12}$~; \quad
$\dfrac{3^2 \times 3^7}{3^5} = 3^{9-5} = 3^4$~; \quad
$2^5 \times 4 = 2^5 \times 2^2 = 2^7$.
\end{solution}

\begin{solution}{exo:g9:fractions:5}
$45\,000 = 4.5 \times 10^4$~; \quad
$0.0072 = 7.2 \times 10^{-3}$~; \quad
$302.5 = 3.025 \times 10^2$~; \quad
$0.000\,000\,9 = 9 \times 10^{-7}$~; \quad
douze millions $= 1.2 \times 10^7$.
\end{solution}

\begin{solution}{exo:g9:fractions:6}
$(2 \times 10^7) \times (4.5 \times 10^{-3}) = 9 \times 10^{4}$.

$\dfrac{9 \times 10^5}{3 \times 10^{-2}} = 3 \times 10^{5 - (-2)}
= 3 \times 10^7$.

$(5 \times 10^3)^2 = 25 \times 10^6 = 2.5 \times 10^7$.
\end{solution}

\begin{solution}{exo:g9:fractions:7}
L'eau ajoutée représente
$\dfrac34 - \dfrac35 = \dfrac{15}{20} - \dfrac{12}{20} = \dfrac{3}{20}$
du réservoir. Donc $\frac{3}{20}$ de la capacité vaut $30$ litres~:
$\frac{1}{20}$ vaut $10$ litres, et la capacité est $200$ litres.
\end{solution}

\begin{solution}{exo:g9:fractions:8}
Après Alice~: $1 - \frac14 = \frac34$ du gâteau \omterm{def:g3:division:remainder}{reste}. Ben mange
$\frac13 \times \frac34 = \frac14$, laissant
$\frac34 - \frac14 = \frac12$. Carla mange la \omterm{def:g3:division:half}{moitié} de cela,
$\frac14$, laissant $\frac14$ du gâteau.
\end{solution}

\begin{solution}{exo:g9:fractions:9}
Temps $= \dfrac{\text{distance}}{\text{vitesse}} =
\dfrac{1.5 \times 10^{11}}{3 \times 10^8} = 0.5 \times 10^3 = 500$
secondes, soit $8$ minutes et $20$ secondes.
\end{solution}

\begin{solution}{exo:g9:fractions:10}
Écrire les deux nombres comme \omterm{def:g9:fractions:power}{puissances} d'\omterm{def:g8:powers:def}{exposant} $10$~:
\[
2^{30} = \left(2^3\right)^{10} = 8^{10},
\qquad
3^{20} = \left(3^2\right)^{10} = 9^{10}.
\]
Puisque $8 < 9$, en multipliant dix copies de chacun~: $8^{10} <
9^{10}$, donc $2^{30} < 3^{20}$.
\end{solution}

\begin{solution}{pb:g9:fractions:1}
\textbf{1.} $\frac14 = 0.25$ (s'arrête)~; $\frac13 = 0.333\ldots$
(période $3$)~; $\frac56 = 0.8333\ldots$ (période $3$ après le
$8$)~; $\frac{3}{11} = 0.272727\ldots$ (période $27$).

\textbf{2.} À chaque étape de la division par $b$, le \omterm{def:g3:division:remainder}{reste} est
l'une des $b$ valeurs $0, 1, \dots, b - 1$
(\cref{thm:g6:wholes:division}). Si un \omterm{def:g3:division:remainder}{reste} $0$ apparaît, la
division s'arrête. Sinon, après au plus $b$ étapes, un \omterm{def:g3:division:remainder}{reste} doit
réapparaître --- il y a plus d'étapes que de \omterm{def:g3:division:remainder}{restes} possibles --- et
à partir de cet instant la division reproduit exactement la suite de
chiffres qu'elle avait produite après la première apparition~:
l'écriture décimale tourne en boucle indéfiniment. S'arrêter ou se
répéter~: il n'existe pas de troisième comportement.

\textbf{3.} $\frac17 = 0.142857\,142857\ldots$~: la période
$142857$ est de longueur $6$. La question 2 promettait une
répétition en au plus $7$ étapes --- ici les \omterm{def:g3:division:remainder}{restes}
$1, 3, 2, 6, 4, 5$ apparaissent tous avant que le $1$ ne revienne~:
la plus longue mélopée possible pour une division par $7$.

\textbf{4.} $\frac ab$ s'arrête exactement lorsque l'une des
\omterm{def:g9:fractions:fraction}{fractions} $\frac{a \times k}{b \times k}$ a pour \omterm{def:g4:fractions:def}{dénominateur} $10$,
$100$, $1000, \dots$ --- c'est-à-dire lorsque le \omterm{def:g4:fractions:def}{dénominateur} peut
être multiplié jusqu'à une \omterm{def:g9:fractions:power}{puissance} de dix, comme
$4 \times 25 = 100$ ou $8 \times 125 = 1\,000$. Pour $3$ et $6$,
c'est sans espoir~: une \omterm{def:g9:fractions:power}{puissance} de dix a pour \omterm{def:g1:addition:def}{somme} des chiffres
$1$, alors que tout \omterm{def:g5:division:multiple}{multiple} de $3$ a une \omterm{def:g1:addition:def}{somme} de chiffres \omterm{def:g5:division:multiple}{multiple}
de $3$ --- aucun \omterm{def:g5:division:multiple}{multiple} de $3$ (donc aucun de $6$) ne vaut jamais
$10$, $100$, $1\,000, \dots$ Pour $11$~: diviser $10$, $100$,
$1\,000$, $10\,000$ par $11$ laisse toujours un \omterm{def:g3:division:remainder}{reste} ($10$, $1$,
$10$, $1, \dots$), jamais $0$. Les \omterm{def:g4:fractions:def}{dénominateurs} qui atteignent une
\omterm{def:g9:fractions:power}{puissance} de dix s'arrêtent~; tous les autres se répètent.

\textbf{5.} $\frac{9}{40}$~: $40 \times 25 = 1\,000$, s'arrête~:
$\frac{225}{1000} = 0.225$. $\frac{33}{50}$~:
$50 \times 2 = 100$, s'arrête~: $\frac{66}{100} = 0.66$.
$\frac{5}{12}$~: $12$ est un \omterm{def:g5:division:multiple}{multiple} de $3$, se répète
($0.41666\ldots$).

\textbf{6.} $10x = 7.777\ldots$, donc $10x - x = 7$ (les queues se
compensent parfaitement), $9x = 7$ et $x = \frac79$. Vérification
par la division~: $7 \div 9 = 0.777\ldots$

\textbf{7.} La période a deux chiffres~: on décale donc de $100$~:
$100x - x = 72.7272\ldots - 0.7272\ldots = 72$, d'où
$99x = 72$ et $x = \frac{72}{99} = \frac{8}{11}$.

\textbf{8.} $1000x = 583.333\ldots$ et $100x = 58.333\ldots$~; en
soustrayant, $900x = 525$, donc
$x = \frac{525}{900} = \frac{7}{12}$ (diviser par $75$). On a bien
$\frac{7}{12} = 0.58333\ldots$

\textbf{9.} $10x - x = 9.999\ldots - 0.999\ldots = 9$, donc
$9x = 9$ et $x = 1$. Ce n'est pas «~juste en dessous~» de $1$~:
$0.999\ldots$ \emph{est} le nombre $1$, écrit dans un second
costume. Tout nombre strictement inférieur à $1$ laisse un écart, et
\cref{pb:g6:decimals:1} a montré qu'un écart contient toujours
d'autres nombres --- alors que rien ne tient entre $0.999\ldots$ et
$1$. Les neufs sans fin sont le \omterm{def:g3:division:remainder}{reste} de chocolat du
\cref{pb:g6:fractions:1} rétréci en dessous de toute quantité
positive~: il ne reste rien.

\textbf{10.} $2.4545\ldots = 2 + \frac{45}{99} = 2 +
\frac{5}{11} = \frac{27}{11}$. Et
$0.123123\ldots = \frac{123}{999} = \frac{41}{333}$ (diviser par
$3$).

\textbf{11.} Dictionnaire~: $0.037037\ldots = \frac{37}{999}$.
Comme $999 = 27 \times 37$, cela se simplifie en $\frac{1}{27}$~:
c'est la \omterm{def:g9:fractions:fraction}{fraction} $\frac{1}{27}$ qui psalmodie «~$037$~».
Vérification par la division~: $1 \div 27 = 0.037037\ldots$

\textbf{12.} $0.0272727\ldots = \frac{1}{10} \times
0.272727\ldots = \frac{1}{10} \times \frac{27}{99}
= \frac{27}{990} = \frac{3}{110}$ (diviser par $9$).

\textbf{13.} Les blocs de \omterm{def:g1:counting:zero}{zéros} s'allongent sans fin~: aucune
période fixe ne peut donc se répéter indéfiniment. Ce n'est
\emph{pas} un décimal périodique (et il ne s'arrête jamais).
D'après la partie I, toute \omterm{def:g9:fractions:fraction}{fraction} s'arrête ou se répète --- ce
nombre n'est donc pas une \omterm{def:g9:fractions:fraction}{fraction} du tout~: il est irrationnel,
fabriqué sur mesure. La vedette de l'irrationalité, $\sqrt2$, est
démasquée au \cref{ch:g9:sqrt}.

\textbf{14.} $2^{30} = \left(2^{10}\right)^3 \approx
(1.024)^3 \times 10^9 \approx 1.07 \times 10^9$~: un nombre
légèrement supérieur à $10^9$, il a donc $10$ chiffres.
(Exactement~: $2^{30} = 1\,073\,741\,824$.)

\textbf{15.} $0.20262026\ldots = \frac{2026}{9999}$, déjà
irréductible ($9999 = 3 \times 3 \times 11 \times 101$ ne partage
aucun facteur avec $2026 = 2 \times 1013$). Le dictionnaire
achevé~: \emph{les \omterm{def:g9:fractions:fraction}{fractions} sont exactement les décimaux qui
s'arrêtent ou se répètent} --- les \omterm{def:g4:fractions:def}{dénominateurs} amis des \omterm{def:g9:fractions:power}{puissances}
de dix s'arrêtent, tous les autres psalmodient une période~; et
réciproquement, toute mélopée, même celle d'un millésime, est une
\omterm{def:g9:fractions:fraction}{fraction} déguisée.
\end{solution}
